% ============================================================================
% PLIK NIEAKTYWNY — rozdział 5 został usunięty z pracy decyzją usera (2026-05-28).
% Powód: 3 case studies bez statystyki nie stanowiły rzeczywistej walidacji,
% a poszczególne treści dublują się z innymi rozdziałami:
%   - Lista 12 współczynników + geneza → rozdz. 3.5
%   - Walidacja literaturowa → rozdz. 3.5 (wartości referencyjne)
%   - Case study Film 10 → rozdz. 4 (aspect ratio fix)
%   - Janek edge case → rozdz. 3.6 (reguła łączona SI+confidence)
%   - Ograniczenia obliczeń → rozdz. 8 (dyskusja)
%   - Pipeline pod różnymi warunkami → rozdz. 7 (wrażliwość)
%
% Plik zachowany jako draft do ewentualnej kanibalizacji fragmentów.
% Nie jest inputowany w main.tex.
% ============================================================================

\chapter{Współczynniki biegu --- wyniki i~walidacja}
\label{ch:wspolczynniki}

\section{Wprowadzenie}
\label{sec:wsp-wprowadzenie}

Niniejszy rozdział odpowiada na pierwsze pytanie badawcze postawione we wstępie: jak zaprojektować pipeline maszynowego uczenia, który na podstawie nagrania z~jednej kamery (\textit{monocular 2D}) wylicza zestaw współczynników biomechanicznych opisujących technikę biegu. W~rozdziale~\ref{ch:materialy-metody} (sekcja~\ref{sec:obliczanie-wspolczynnikow}) opisałem metody obliczania każdego współczynnika z~osobna. Tutaj pokazuję, jakie wartości pipeline produkuje w~praktyce, kiedy zostaje uruchomiony na rzeczywistych nagraniach, oraz oceniam ich wiarygodność przez porównanie z~zakresami raportowanymi w~literaturze biomechanicznej.

Pipeline wylicza łącznie 12~współczynników podzielonych na dwie grupy, wraz z~dodatkowymi wskaźnikami symetrii lewo-prawo.

\textbf{Pięć współczynników temporalnych} obliczanych z~sekwencji etykiet faz biegu i~częstotliwości klatek wideo:
\begin{itemize}
    \item kadencja,
    \item czas kontaktu z~podłożem (GCT) --- osobno dla lewej i~prawej nogi,
    \item czas lotu,
    \item czas cyklu (oraz długość kroku, jeśli użytkownik poda prędkość bieżni),
    \item duty factor.
\end{itemize}

\textbf{Siedem współczynników przestrzennych} obliczanych z~geometrii keypointów MediaPipe Pose:
\begin{itemize}
    \item trzy kąty stawów (kolano, biodro, kostka) --- na każdą nogę osobno, uśrednione w~obrębie fazy,
    \item kąt kolana w~momencie kontaktu z~podłożem (\textit{initial contact}),
    \item pochylenie tułowia,
    \item oscylacja pionowa bioder, znormalizowana długością tułowia,
    \item wzorzec lądowania stopy (\textit{heel} / \textit{midfoot} / \textit{forefoot strike}).
\end{itemize}

Dodatkowo dla każdego współczynnika obliczanego osobno dla lewej i~prawej strony ciała wyznaczam wskaźnik symetrii Robinsona (SI), opisany w~sekcji~\ref{sec:symmetry}.

Pierwotnie planowałem dużo węższy zakres --- cztery do pięciu podstawowych metryk (kadencja, GCT, długość kroku, wzorzec lądowania stopy). W~trakcie pracy nad pipeline okazało się jednak, że kolejne współczynniki dają się obliczyć w~naturalny sposób z~tych samych danych wejściowych (sekwencja faz + wygładzone keypointy), bez dodatkowych kosztów obliczeniowych ani dodatkowych wymagań wobec nagrania. Rozszerzony zestaw 12~współczynników daje pełniejszy obraz techniki biegu i~pozwala silnikowi reguł rekomendacji (rozdział~\ref{ch:rekomendacje}) wykrywać szerszy zakres potencjalnych problemów.

\section{Strategia walidacji}
\label{sec:strategia-walidacji}

Najbardziej miarodajną metodą walidacji byłoby porównanie wartości wyliczonych z~monocular 2D wideo z~pomiarami trójwymiarowego systemu motion capture (3D MoCap), takiego jak Vicon lub Qualisys, traktowanymi jako ground truth. Dostęp do takiego laboratorium wymaga znacznego budżetu i~specjalistycznej infrastruktury, co wykracza poza możliwości pracy magisterskiej realizowanej w~warunkach domowych. Z~tego samego powodu w~zbiorze nie ma żadnego nagrania, dla którego istniałyby równoległe pomiary 3D, więc nie mogę raportować precyzji w~jednostkach fizycznych (np.~MAE w~milisekundach lub stopniach).

Zastosowałem walidację pośrednią: dla każdego współczynnika sprawdzam, czy wyliczona wartość mieści się w~zakresie publikowanym w~literaturze biomechanicznej --- przede wszystkim Novacheck~\cite{novacheck1998}, Heiderscheit~\cite{heiderscheit2011}, Souza~\cite{souza2016} i~Daoud~\cite{daoud2012}. Zakresy te są zebrane w~tabeli referencyjnej (sekcja~\ref{sec:obliczanie-wspolczynnikow}) i~służą zarówno jako kryterium oceny, jak i~jako progi dla silnika reguł rekomendacji.

Walidację jakościową przeprowadziłem na trzech biegaczach o~świadomie zróżnicowanej technice: zaawansowanym (Adam, zbiór treningowy), pośrednim (film~10 ze~zbioru testowego) i~z~wyraźnymi problemami (Janek, przypadek brzegowy spoza zbioru). Celem nie było dowodzenie pełnej precyzji pipeline'u, lecz sprawdzenie, czy potrafi on \textit{różnicować} biegaczy --- czyli czy lepsza technika prowadzi do wartości bliższych zakresom z~literatury, a~problemy techniczne wywołują sygnały ostrzegawcze.

\section{Studium przypadków --- trzech biegaczy}
\label{sec:case-studies}

Tabela zbiorcza wszystkich wartości znajduje się w~sekcji~\ref{sec:porownanie-zbiorcze}. Tutaj przedstawiam każdy przypadek osobno, koncentrując się na wnioskach dotyczących pipeline'u, a~nie na pełnej charakterystyce techniki biegu konkretnego biegacza.

\subsection{Adam --- biegacz zaawansowany}
\label{sec:case-adam}

Pierwszym biegaczem walidacyjnym jest Adam, członek klubu biegowego Zabiegani na Szamę. Jego nagranie (film~12, 80~s, FPS~$\approx$~30) trafiło do zbioru treningowego, więc dla wyników klasyfikatora faz nie traktuję go jako oceny generalizacji. W~kontekście walidacji współczynników biomechanicznych jest jednak przydatne jako sprawdzian, czy pipeline produkuje sensowne wartości dla biegacza zaawansowanego.

Wszystkie wskaźniki temporalne mieszczą się w~zakresach typowych dla biegacza zaawansowanego: kadencja 173~spm (literatura: 170--185), GCT lewej nogi 203~ms i~prawej 245~ms (zakres rekreacyjnego i~szybkiego biegu, 160--280~ms), oscylacja pionowa znormalizowana 0{,}14 długości tułowia (dobry biegacz: 0{,}12--0{,}16). Foot strike na obu nogach to forefoot, z~wyraźną dominacją w~96\% kontaktów lewej i~75\% prawej. Kąt kolana w~momencie kontaktu (161° lewa, 162° prawa) mieści się w~zalecanym zakresie 160--175°.

Jedna obserwacja jest jednak nietrywialna: wskaźnik symetrii GCT wynosi SI~=~18{,}7\%, co przekracza próg 10\% i~zostało oznaczone jako potencjalny problem. Pipeline samodzielnie nie rozstrzyga, czy biegacz rzeczywiście ma asymetryczne kontakty, czy jest to artefakt klasyfikatora przesuwającego granice faz o~1--2~klatki. Analiza wrażliwości (rozdział~\ref{ch:wrazliwosc}) pokazuje, że w~tym wypadku różnica wynika w~znacznej mierze z~drugiego źródła. Przykład ten dobrze ilustruje, że walidacja przez zakresy literatury mówi nam, \textit{kiedy} wartość wykracza poza normę, ale nie zawsze \textit{dlaczego}.

\subsection{Film 10 --- analiza z~fizjoterapeutą}
\label{sec:case-film10}

Drugim biegaczem walidacyjnym jest osoba z~filmu~10 (zbiór testowy, $\approx$~320~klatek, FPS~$\approx$~30, ujęcie pionowe $608 \times 1080$~px). Wybrałem ten przypadek celowo, ponieważ orientacja kamery uwypukla wrażliwość różnych grup współczynników na warunki akwizycji.

Wartości temporalne pozostają sensowne: kadencja 163~spm (zakres rekreacyjny), GCT 285/226~ms (jogging do szybkiego biegu), czas lotu 116~ms, pochylenie tułowia 9{,}9°. Pipeline temporalny radzi sobie z~filmem pionowym, ponieważ czasy faz nie zależą od orientacji kadru.

Współczynniki przestrzenne pokazują jednak limity podejścia 2D. Kąt kolana w~momencie kontaktu wynosi 98° (lewa) i~137° (prawa) --- obie wartości daleko poza zakresem prawidłowym 160--175° i~oznaczone jako problem ``siedzącego biegu''. Realnie wynika to z~połączenia perspektywy kamery z~techniką biegacza: w~ujęciu pionowym geometria kątów na płaszczyźnie obrazu jest zniekształcona względem rzeczywistej pozycji nogi w~przestrzeni. Najbardziej spektakularne są kąty stopy w~detekcji wzorca lądowania --- $-97^\circ$ i~$-98^\circ$ w~wartości bezwzględnej, przekraczające próg wiarygodności $45^\circ$ wielokrotnie. Pipeline poprawnie zapala flagę \texttt{low\_confidence} (sekcja~\ref{sec:obliczanie-wspolczynnikow}) i~oznacza wzorzec lądowania jako niewiarygodny. Ten przypadek pokazuje istotną cechę projektową: pipeline jest świadomy własnych ograniczeń i~nie udaje precyzji, której fizycznie nie może zapewnić.

\subsection{Janek --- przypadek brzegowy}
\label{sec:case-janek}

Trzeci biegacz, Janek, jest przypadkiem szczególnym. Jego nagranie (FPS~30, $\approx$~2400~klatek) nie znalazło się w~żadnym ze~zbiorów --- ani treningowym, ani walidacyjnym, ani testowym --- ponieważ posłużyło wyłącznie do oceny pipeline'u jako całości na nieznanym wcześniej materiale.

Wyniki pipeline'u są wyraźnie odmienne od dwóch poprzednich przypadków: kadencja 148~spm (poniżej progu 150 oznaczonego jako bardzo niska, sygnalizująca overstriding lub chód), GCT lewej nogi 192~ms wobec prawej 357~ms (różnica 165~ms), duty factor lewa 0{,}24 wobec prawej 0{,}44, kąt kolana w~kontakcie 138° lewa i~148° prawa (oba poniżej zakresu 160--175°). Wskaźnik symetrii Robinsona dla GCT wynosi SI~=~59{,}8\%, dla duty factor 60{,}0\%. Wartości tego rzędu nie są spotykane u~zdrowych biegaczy --- to różnice z~obszaru patologii ruchu albo poważnego błędu pomiaru.

Pipeline poprawnie zareagował: silnik reguł rekomendacji wygenerował flagę \texttt{critical} dla asymetrii kroków oraz dla łącznej kombinacji wysokiego SI z~obniżoną pewnością predykcji (\texttt{avg\_confidence}~=~0{,}88), korzystając z~reguły łączonej opisanej w~sekcji~\ref{sec:rules-combined}. Ta sygnalizacja jest tu właściwa niezależnie od tego, czy źródłem są błędy klasyfikatora na nietypowym nagraniu, czy realny problem z~techniką biegu --- w~obu przypadkach wynik wymaga interwencji człowieka i~nie powinien być prezentowany użytkownikowi jako wiarygodna analiza biomechaniczna.

\section{Porównanie zbiorcze}
\label{sec:porownanie-zbiorcze}

Tabela~\ref{tab:wsp-zbiorcza} zestawia kluczowe współczynniki dla trzech biegaczy walidacyjnych.

\begin{table}[htbp]
\centering
\caption{Wyniki obliczeń współczynników dla trzech biegaczy walidacyjnych}
\label{tab:wsp-zbiorcza}
\small
\begin{tabular}{lccc}
\hline
\textbf{Współczynnik} & \textbf{Adam} & \textbf{Film~10} & \textbf{Janek} \\
\hline
Kadencja [spm]                       & 173      & 163      & 148 (🔴) \\
GCT lewa [ms]                        & 203      & 285      & 192 \\
GCT prawa [ms]                       & 245      & 226      & 357 (🔴) \\
Czas lotu [ms]                       & 122      & 116      & 130 \\
Duty factor L / P                    & 0{,}30 / 0{,}36 & 0{,}39 / 0{,}31 & 0{,}24 / 0{,}44 \\
Kąt kolana @ kontakt L [$^\circ$]    & 162      & 98 (🟡)  & 138 (🟡) \\
Kąt kolana @ kontakt P [$^\circ$]    & 162      & 137 (🟡) & 149 (🟡) \\
Pochylenie tułowia [$^\circ$]        & 2{,}3 (🟡) & 9{,}9    & 9{,}9 \\
Vertical osc. (per torso)            & 0{,}14   & 0{,}10 (🟡) & 0{,}11 (🟡) \\
Foot strike L / P                    & FF / FF  & FF / FF$^\dagger$ & MF / MF \\
SI GCT [\%]                          & 18{,}7   & 22{,}8   & 59{,}9 (🔴) \\
SI duty factor [\%]                  & 18{,}7   & 21{,}4   & 60{,}0 (🔴) \\
Flagi krytyczne (\texttt{critical})  & 0        & 0        & 2 \\
\hline
\multicolumn{4}{l}{\footnotesize FF = forefoot, MF = midfoot. $^\dagger$ Flaga \texttt{low\_confidence} (kąt stopy $>$~$45^\circ$ b.w.).} \\
\end{tabular}
\end{table}

Z~zestawienia wynikają trzy obserwacje. Po pierwsze, pipeline różnicuje trzech biegaczy: każda kolumna ma inną sygnaturę odchyleń od literatury i~inną liczbę zapalonych flag krytycznych. Po drugie, wszystkie wartości temporalne mieszczą się w~zakresach raportowanych w~literaturze i~odpowiadają intuicji biomechanicznej --- Adam ma najwyższą kadencję i~najkrótsze kontakty, Janek najniższą kadencję i~najbardziej asymetryczne kontakty. Po trzecie, współczynniki przestrzenne (zwłaszcza kąty kolana w~momencie kontaktu i~wzorzec lądowania stopy) okazują się wyraźnie wrażliwe na warunki nagrania --- u~biegacza z~filmu~10 dają wartości fizycznie nieprawdopodobne, ale pipeline rozpoznaje to przez flagę \texttt{low\_confidence}.

Wskaźnik symetrii Robinsona poprawnie wskazuje skalę problemu: 18--22\% u~biegaczy bez większych zastrzeżeń biomechanicznych i~60\% w~przypadku Janka, gdzie ani biegacz, ani warunki nagrania nie są typowe.

\section{Ograniczenia obliczeń}
\label{sec:ograniczenia-obliczen}

\begin{itemize}
    \item \textbf{Brak kalibracji piksel--metr.} Oscylacja pionowa raportowana jest w~jednostkach znormalizowanych długością tułowia, a~nie w~centymetrach. Przeliczenie na centymetry wymagałoby znajomości wzrostu biegacza i~odległości kamery od niego, co rzadko bywa dostępne w~praktycznym scenariuszu użycia.
    \item \textbf{Monocular 2D --- artefakty perspektywy.} Wszystkie współczynniki kątowe wyliczane są w~płaszczyźnie obrazu, nie w~przestrzeni trójwymiarowej. Przy ujęciu innym niż ściśle boczne (np.~wideo pionowe z~filmu~10) geometria kątów ulega zniekształceniu. Pipeline kompensuje to częściowo przez flagę \texttt{low\_confidence} dla kąta stopy, ale nie dla kątów stawów.
    \item \textbf{Stride length wymaga inputu użytkownika.} Na bieżni mechanicznej biegacz porusza się w~miejscu, więc długości kroku nie da się wyliczyć z~samego wideo --- pipeline wymaga podania prędkości bieżni przez użytkownika. Dla nagrań bez tej informacji długość kroku nie jest raportowana.
    \item \textbf{Brak ground truth dla precyzji.} Walidacja przez zakresy literatury pozwala stwierdzić, czy wartość jest realistyczna biomechanicznie, ale nie pozwala podać dokładności rzędu MAE w~jednostkach fizycznych. Dla pełnej oceny wymagana byłaby równoległa rejestracja 3D MoCap, co pozostawiam jako kierunek dalszych prac.
    \item \textbf{Non-determinizm MediaPipe.} Domyślny silnik (XNNPACK delegate) wprowadza drobne, nieprzewidywalne różnice między uruchomieniami dla tego samego nagrania. Przy 30~FPS i~$\approx$~12~klatkach na cykl różnice te są zwykle pomijalne, ale w~filmach z~niskim FPS mogą się kumulować.
\end{itemize}
